
\section{Why a Church that has received the truth still has to teach}

Christian faith is not assent to a classification system. It is the graced
answer of the whole person to God who reveals himself in Jesus Christ. Yet a
faith confessed in history needs words, judgments, teachers, and practices.
The apostles preached a determinate Gospel; churches guarded a rule of faith;
bishops judged rival accounts of Christ; councils chose language capable of
excluding an error without pretending to exhaust the mystery. The resulting
distinctions do not divide truth into important and disposable parts. They say
what kind of claim is being made, by whom, with what relation to Revelation,
and what response the claim therefore asks of a Catholic.

The distinction serves natural human goods. It makes speech intelligible,
preserves institutional memory, permits real disagreement to be located,
protects consciences from inflated claims, and prevents every policy or
theologian's conclusion from masquerading as divine speech. These goods do not
produce Revelation. They are created capacities---reason, language, testimony,
authority, and common life---placed at the service of receiving it.

Its supernatural purpose is deeper. Christ entrusted the apostolic deposit to
the Church so that the Gospel would be preached faithfully, the sacraments
celebrated, believers formed in one confession, error corrected, and persons
led to communion with God. The Magisterium is therefore ministerial: \emph{Dei
Verbum} 10 says it is not above the word of God but serves it. Assistance by
the Holy Spirit is not a replacement for evidence, deliberation, holiness, or
responsible language; it is Christ's gift for the Church's fidelity under the
conditions he has willed.

